\documentclass[11pt]{article}
\usepackage[a4paper,margin=1in]{geometry}
\usepackage[T1]{fontenc}
\usepackage{lmodern}
\usepackage{amsmath,amssymb,amsthm,booktabs,enumitem,mathtools,microtype}
\usepackage[numbers,sort&compress]{natbib}
\usepackage{xurl}
\usepackage[hidelinks]{hyperref}
\setlength{\emergencystretch}{3em}

\newtheorem{theorem}{Theorem}[section]
\newtheorem{proposition}[theorem]{Proposition}
\newtheorem{lemma}[theorem]{Lemma}
\newtheorem{corollary}[theorem]{Corollary}
\theoremstyle{definition}
\newtheorem{definition}[theorem]{Definition}
\theoremstyle{remark}
\newtheorem{remark}[theorem]{Remark}

\newcommand{\mob}{\mu}
\newcommand{\Thetaqr}{\Theta_{q,r}}
\newcommand{\kap}{\kappa}
\newcommand{\ind}{\mathbf 1}
\newcommand{\Zq}{\mathbb Z/q\mathbb Z}
\newcommand{\T}{\{-1,0,+1\}}
\DeclareMathOperator{\lcm}{lcm}

\hypersetup{
  pdftitle={Two-Odd-Factor Terminal-Log Mobius Compiler and the Multi-Shift Squarefree Landscape},
  pdfauthor={RH research program},
  pdfsubject={Terminal logarithmic Mobius diagonalization with at most two odd factors},
  pdfkeywords={Mobius function, terminal logarithmic average, squarefree patterns, finite CRT}
}

\title{Two-Odd-Factor Terminal-Log M\"obius Compiler and the Multi-Shift Squarefree Landscape}
\author{RH research program}
\date{August 10, 2026}

\begin{document}
\maketitle

\begin{abstract}
For any fixed finite set of distinct integer shifts, we diagonalize terminal
logarithmic averages of coordinatewise-quadratic M\"obius polynomials whose
monomials contain at most two odd exponents.  The total polynomial degree may
be as large as twice the number of shifts.  All channels with one or two odd
M\"obius factors vanish; the surviving all-even channels are exact
multi-shift squarefree phase densities.  Their Euler factors retain the
collision distinction between congruence classes modulo \(p\) and modulo
\(p^2\).

The proof combines a local finite-CRT calculation with a uniform Boolean
square-mask tail.  After fixing the cutoff, one odd factor is handled by the
frozen periodic one-form cancellation in RH-392 and two odd factors by its
arbitrary-nonzero-determinant terminal theorem.  The order of limits is
cutoff fixed, terminal limit, and then cutoff removal.  In three variables
the compiler contains \(26\) of the \(27\) ternary interpolation monomials;
the sole excluded coefficient is the signed-cube coefficient \(c_{111}\).
This gives a concrete cancellation class of exactly \(192\) among the
\(512\) two-input truth tables used in a distinguished-current score.

We also determine the global squarefree-density landscape.  For at most
three shifts it has a positive sharp lower product, with equality exactly
when all pairwise differences are squarefree.  From four shifts onward zero
is attained precisely through a complete residue cover modulo some
\(p^2\).  For every fixed number of at least two shifts the supremum is
\(6/\pi^2\), approached but never attained.  Every datum is fixed before the
terminal limit; no higher odd correlation, growing family, ordinary
Ces\`aro, or multishift safe-table capacity is claimed.
\end{abstract}

\noindent\textbf{Keywords:} M\"obius function; terminal logarithmic average;
squarefree patterns; finite Chinese remainder theorem; ternary interpolation.

\section{Definitions and main results}

Write \(\mob_0(t)=\mob(t)\) for integers \(t\ge1\), and
\(\mob_0(t)=0\) for \(t\le0\).  A terminal clock is a function satisfying
\begin{equation}
 1\le \omega(X)\le X,\qquad \omega(X)\longrightarrow\infty.
 \label{eq:clock}
\end{equation}
Fix integers \(m,q\ge1\), pairwise-distinct integer shifts
\(a_1,\ldots,a_m\), and fixed \(q\)-periodic coefficients
\(c_\alpha(r)\).  All these data, including the clock, are fixed before
\(X\to\infty\).  Put
\[
 z_i(n)=\mob_0(n-a_i),\qquad
 O(\alpha)=\{i:\alpha_i=1\},\qquad
 E(\alpha)=\{i:\alpha_i=2\}.
\]
For polynomials
\begin{equation}
 P_r(z)=\sum_{\substack{\alpha\in\{0,1,2\}^m\\|O(\alpha)|\le2}}
 c_\alpha(r)\prod_{i=1}^m z_i^{\alpha_i},
 \label{eq:poly}
\end{equation}
define the terminal functional
\begin{equation}
 T_X(P;\omega)=\frac1{\log\omega(X)}
 \sum_{X/\omega(X)<n\le X}\frac{P_{n\bmod q}(z_1(n),\ldots,z_m(n))}{n}.
 \label{eq:functional}
\end{equation}

For \(E\subseteq\{1,\ldots,m\}\), let
\begin{equation}
 B_p(E)=\{a_i\bmod p^2:i\in E\},\quad
 \nu_p(E)=|B_p(E)|,
 \quad \tau_{p,E}(r)=|\{b\in B_p(E):b\equiv r\pmod p\}|.
 \label{eq:local-data}
\end{equation}
The braces in \(B_p(E)\) are essential: residues are first deduplicated
modulo \(p^2\), while two distinct such residues may still collide modulo
\(p\).  Define
\begin{align}
 \Theta_{q,r}(E)
  &=\frac1q\prod_{p\nmid q}\left(1-\frac{\nu_p(E)}{p^2}\right)
    \prod_{p\parallel q}\left(1-\frac{\tau_{p,E}(r)}p\right)
    \prod_{p^2\mid q}\ind_{\{r\bmod p^2\notin B_p(E)\}},
 \label{eq:theta}\\
 \kap_E&=\prod_p\left(1-\frac{\nu_p(E)}{p^2}\right).
 \label{eq:kappa}
\end{align}
Thus \(\Theta_{q,r}(\varnothing)=1/q\).

\begin{theorem}[Two-odd-factor terminal compiler]
\label{thm:compiler}
For every fixed datum above and every clock \eqref{eq:clock},
\begin{equation}
 T_X(P;\omega)\longrightarrow
 \sum_{r\bmod q}\ \sum_{\alpha\in\{0,2\}^m}
 c_\alpha(r)\Theta_{q,r}(E(\alpha)).
 \label{eq:compiler}
\end{equation}
Every admissible monomial with one or two odd exponents vanishes.  The
all-even channels are the only survivors, and
\begin{equation}
 \sum_{r\bmod q}\Theta_{q,r}(E)=\kap_E.
 \label{eq:phase-sum}
\end{equation}
\end{theorem}

The number of admitted monomials is
\begin{equation}
 D_m=2^m+m2^{m-1}+\binom m2 2^{m-2},
 \label{eq:dimension}
\end{equation}
where terms with impossible support sizes are interpreted as zero.  In
particular, \(D_3=26\): only \(z_1z_2z_3\) is absent.

\begin{corollary}[The signed-cube criterion]
\label{cor:c111}
Let \(f:\T^3\to\mathbb R\), and let its unique coordinatewise-quadratic
interpolant have coefficient \(c_{111}(f)\).  Then
\begin{equation}
 c_{111}(f)=\frac18\sum_{\varepsilon\in\{-1,+1\}^3}
 \varepsilon_1\varepsilon_2\varepsilon_3f(\varepsilon).
 \label{eq:c111}
\end{equation}
The interpolant is covered by Theorem~\ref{thm:compiler} exactly when
\(c_{111}(f)=0\).
\end{corollary}

\begin{corollary}[A \(192\)-table distinguished-current class]
\label{cor:192}
Fix three distinct shifts and, in every phase \(r\), a table
\(f_r:\T^2\to\{-1,+1\}\) satisfying
\begin{equation}
 f_r(1,1)-f_r(1,-1)-f_r(-1,1)+f_r(-1,-1)=0.
 \label{eq:corner}
\end{equation}
Then, for every terminal clock,
\begin{equation}
 \frac1{\log\omega(X)}\sum_{X/\omega(X)<n\le X}
 \frac{\mob_0(n-a_3)
 f_{n\bmod q}(\mob_0(n-a_1),\mob_0(n-a_2))}{n}
 \longrightarrow0.
 \label{eq:table-limit}
\end{equation}
Exactly \(192\) of the \(512\) tables satisfy \eqref{eq:corner}.  The other
\(320\) tables are outside this theorem; no assertion about their limits is
made.
\end{corollary}

For a finite set \(A=\{a_1,\ldots,a_m\}\) of distinct integers, write
\begin{equation}
 \kap_A=\prod_p\left(1-\frac{|A\bmod p^2|}{p^2}\right).
 \label{eq:global-kappa}
\end{equation}

\begin{theorem}[Multi-shift squarefree landscape]
\label{thm:landscape}
The following statements hold.
\begin{enumerate}[label=\textup{(\roman*)}]
 \item If \(1\le m\le3\), then
 \[
  \kap_A\ge C_m:=\prod_p(1-mp^{-2})>0.
 \]
 Equality holds exactly when every nonzero pairwise difference in \(A\) is
 squarefree.  The set \(\{0,\ldots,m-1\}\) is a witness.
 \item If \(m\ge4\), then \(\inf_A\kap_A=0\), and zero is attained.
 More precisely, \(\kap_A=0\) exactly when \(A\bmod p^2\) covers all
 \(p^2\) residues for some prime \(p\).  The residues \(0,1,2,3\pmod4\),
 extended by arbitrary distinct shifts, give a witness for every \(m\ge4\).
 \item For \(m=1\), \(\kap_A=6/\pi^2\).  For each fixed \(m\ge2\),
 \[
  \sup_A\kap_A=\frac6{\pi^2},
 \]
 but the supremum is not attained.
\end{enumerate}
An individual phase density \(\Theta_{q,r}(E)\) may vanish even when
\(\kap_E>0\).
\end{theorem}

\section{Squarefree phase densities}
\label{sec:density}

\begin{lemma}[Finite-prime phase count]
\label{lem:finite-crt}
Let \(\mathcal P\) be a finite set of primes containing every prime divisor
of \(q\), and put \(M=\lcm(q,\prod_{p\in\mathcal P}p^2)\).  The density in
the phase \(n\equiv r\pmod q\) of integers avoiding
\(n\equiv a_i\pmod{p^2}\), for all \(i\in E\) and
\(p\in\mathcal P\), equals the finite truncation of \eqref{eq:theta}.
Summing these phase densities over \(r\pmod q\) gives
\(\prod_{p\in\mathcal P}(1-\nu_p(E)/p^2)\).
\end{lemma}

\begin{proof}
Count directly in one complete block modulo \(M\).  The Chinese remainder
theorem separates the prime-square coordinates.  If \(p\nmid q\), exactly
\(\nu_p(E)\) of the \(p^2\) residues are forbidden.  If
\(p\parallel q\), the fixed residue modulo \(p\) leaves \(p\) lifts, of
which \(\tau_{p,E}(r)\) are forbidden.  If \(p^2\mid q\), the phase is
either forced into a forbidden class or avoids all of them.  These are the
three factors in \eqref{eq:theta}.  Summing over the compatible phase
coordinates restores the unrestricted factor at every prime.
\end{proof}

\begin{lemma}[Cutoff removal]
\label{lem:tail}
For fixed shifts and a cutoff \(P\), the terminal harmonic mass on which at
least one \(n-a_i\), \(i\in E\), is divisible by \(p^2\) for some prime
\(p>P\) is
\begin{equation}
 O_{m,\mathbf a}\!\left(\frac{\log\omega(X)}P+1\right).
 \label{eq:tail}
\end{equation}
The corresponding prefix counting error is
\(O_{m,\mathbf a}(N/P+\sqrt N)\).
\end{lemma}

\begin{proof}
For the prefix bound, sum \(N/p^2+1\) over
\(P<p\le\sqrt{N+\max_i|a_i|}\), and use a union bound over the fixed set
of shifts.  For the terminal bound, split \((X/\omega(X),X]\) into dyadic
blocks.  On a block \((T,2T]\), each progression contributes
\(O(T/p^2+1)\) points, each of weight \(O(1/T)\).  Summation over primes
and then over the dyadic blocks gives \eqref{eq:tail}; blocks meeting the
fixed lower endpoint contribute only \(O_{m,\mathbf a}(1)\).  The implied constants
may depend on the fixed shifts, but not on \(X\) or \(P\).
\end{proof}

\begin{proposition}[Existence and formula for \(\Theta\)]
\label{prop:theta}
The density of integers \(n\equiv r\pmod q\) for which every
\(n-a_i\), \(i\in E\), is squarefree exists and equals
\(\Theta_{q,r}(E)\) in \eqref{eq:theta}.
\end{proposition}

\begin{proof}
Lemma~\ref{lem:finite-crt} gives the result with finitely many local
conditions.  Lemma~\ref{lem:tail} removes the remaining prime squares.
First let \(N\to\infty\) with \(P\) fixed, and then let \(P\to\infty\).
The same argument after summing phases proves \eqref{eq:phase-sum}.  This
is a local proof; Mirsky's classical pattern theorem is historical context,
not an invoked black box \cite{Mirsky1948}.
\end{proof}

\section{Cancellation with square masks}
\label{sec:cancellation}

We use two frozen inputs from RH-392.  Its equation (19), on printed/PDF
page 5, gives the prefix estimate
\(\sum_{n\le N}\mob_0(n-a)\rho(n)=o(N)\) for every fixed periodic
\(\rho\) and fixed integer shift \(a\).  RH-392, Lemma 2.1, transfers this
prefix estimate to every terminal clock, giving
\begin{equation}
 \sum_{X/\omega(X)<n\le X}\frac{\mob_0(n-a)\rho(n)}n
 =o(\log\omega(X)).
 \label{eq:one-form}
\end{equation}
Its Theorem 2.2 on printed/PDF page 3 gives the analogous terminal statement
for the product of two fixed positive-leading affine forms of arbitrary
nonzero determinant
\cite{RH392}.  The latter ultimately uses Tao's fixed nonparallel
two-point logarithmic theorem \cite{Tao2016LogChowla}; the one-form input
rests on Davenport's classical estimate \cite{Davenport1937}.  TPC-137 is
only a determinant-two proof blueprint and is not being promoted to an
arbitrary-determinant theorem \cite{TPC137}.

\begin{lemma}[Masked odd-channel cancellation]
\label{lem:masked}
Fix \(O,E\subseteq\{1,\ldots,m\}\), with \(O\cap E=\varnothing\), and a
fixed periodic \(\rho\).  If \(1\le|O|\le2\), then
\begin{equation}
 \frac1{\log\omega(X)}\sum_{X/\omega(X)<n\le X}
 \frac{\rho(n)\prod_{i\in O}\mob_0(n-a_i)
       \prod_{j\in E}\mob_0(n-a_j)^2}{n}\longrightarrow0.
 \label{eq:masked}
\end{equation}
\end{lemma}

\begin{proof}
Fix \(P\) and replace each squarefree indicator in the even support by the
finite Boolean condition excluding \(p^2\mid n-a_j\) for \(p\le P\).
The product is then periodic, so it may be absorbed into \(\rho\).  If
\(|O|=1\), equation \eqref{eq:one-form} applies.  If \(|O|=2\), the two
forms \(n-a_i\) and \(n-a_j\) have determinant \(a_i-a_j\ne0\), because
the shifts are distinct, and RH-392, Theorem 2.2 applies.  Values with a
nonpositive argument affect only finitely many terms.

Lemma~\ref{lem:tail} bounds the unnormalized replacement error by
\(O_{m,\mathbf a}(\|\rho\|_\infty(\log\omega/P+1))\).  After division by
\(\log\omega\), keep \(P\)
fixed and let \(X\to\infty\), then let \(P\to\infty\).  This proves
\eqref{eq:masked}.  The fixed-data order is indispensable: neither frozen
input supplies a rate uniform in growing shifts, periods, or masks.
\end{proof}

\begin{proof}[Proof of Theorem~\ref{thm:compiler}]
Expand \eqref{eq:poly}.  When \(O(\alpha)=\varnothing\), the monomial is
the indicator that all shifts indexed by \(E(\alpha)\) are squarefree.
Proposition~\ref{prop:theta} and terminal Abel summation give its limit
\(\Theta_{q,r}(E(\alpha))\) in phase \(r\).  When the odd support has size
one or two, Lemma~\ref{lem:masked} gives zero.  There are finitely many
phases and monomials, so the limits may be summed.  Counting the choices of
zero/two exponents outside an odd support of size \(0,1,2\) gives
\eqref{eq:dimension}.
\end{proof}

For completeness, the terminal Abel step used above follows from partial
summation: if \(A(t)=\alpha t+o(t)\) for a bounded sequence, then its
harmonic mean over \((X/\omega(X),X]\), divided by \(\log\omega(X)\),
tends to \(\alpha\).  When the lower endpoint remains bounded, split at a
fixed large threshold; then \(\log X/\log\omega(X)\to1\).

\section{Interpolation and the 192-table class}

\begin{proof}[Proof of Corollary~\ref{cor:c111}]
The one-variable coefficient of \(x\) in quadratic interpolation on
\(\T\) is \((f(1)-f(-1))/2\).  Taking the tensor product in three
coordinates yields \eqref{eq:c111}.  Among the \(27\) exponent vectors in
\(\{0,1,2\}^3\), the only one with three odd entries is \((1,1,1)\).
Consequently its vanishing is necessary and sufficient for the whole
interpolant to lie in the compiler.
\end{proof}

\begin{proof}[Proof of Corollary~\ref{cor:192}]
Put \(g(x,y,z)=z f(x,y)\).  Tensor interpolation gives
\(c_{111}(g)=c_{11}(f)\), and
\[
 c_{11}(f)=\frac14\{f(1,1)-f(1,-1)-f(-1,1)+f(-1,-1)\}.
\]
If \eqref{eq:corner} holds, every monomial of \(g\) has one or two odd
exponents; it therefore vanishes by Theorem~\ref{thm:compiler}.  For the
census, multiply the four corner signs by their alternating weights.  Their
sum is zero exactly when two of the four transformed signs are positive,
giving \(\binom42=6\) corner patterns.  The five noncorner values are free,
so the count is \(6\cdot2^5=192\).
\end{proof}

\section{Proof of the density landscape}

\begin{proof}[Proof of Theorem~\ref{thm:landscape}]
Write \(\nu_p=|A\bmod p^2|\).  If \(m\le3\), then
\(1-\nu_p/p^2\ge1-m/p^2>0\).  Equality in the product holds exactly when
\(\nu_p=m\) for every prime, which is equivalent to no nonzero pairwise
difference being divisible by a prime square.  This proves part (i), and
the displayed witness has only the squarefree differences \(1\) and \(2\).

Because \(\sum_p\nu_p/p^2<\infty\), an Euler product \eqref{eq:global-kappa}
can vanish only through a zero local factor.  Such a factor is precisely a
complete cover of the \(p^2\) residues.  The four residues modulo \(4\)
give the asserted witness, and adding distinct shifts preserves the cover.
This proves part (ii).

Since \(\nu_p\ge1\), one always has
\(\kap_A\le\prod_p(1-p^{-2})=6/\pi^2\).  For \(m\ge2\), choose a nonzero
difference \(d\).  Some prime does not divide \(d\), so \(\nu_p\ge2\) at
that prime and the inequality is strict.  To approach the upper value, set
\[
 Q_y=\prod_{p\le y}p^2,\qquad A_{m,y}=\{jQ_y:0\le j<m\}.
\]
For \(p\le y\), all shifts coincide modulo \(p^2\), so \(\nu_p=1\).
Once every \(p>y\) has \(p^2>m\),
\[
 \prod_{p\le y}(1-p^{-2})\prod_{p>y}(1-mp^{-2})
 \le\kap_{A_{m,y}}\le\frac6{\pi^2}.
\]
The second product on the left tends to one, proving the supremum claim.
For \(m=1\), every \(\nu_p=1\), so equality is attained.  Finally, a
phase indicator in \eqref{eq:theta} can vanish because of a forced residue
even when the global Euler product is positive; for example, this occurs
for shifts \(\{0,2\}\), \(q=2\), in the even phase.
\end{proof}

\section{Provenance, reproducibility, and scope}

The analytic advance here is a finite-mask synthesis, not a new
higher-order M\"obius input.  RH-392 proves the one- and two-odd-factor
terminal cancellations used in Lemma~\ref{lem:masked}; the present paper
adds arbitrarily many fixed even square masks, the full multi-shift phase
formula, the sharp density landscape, and the \(26/27\) and \(192/512\)
compiler consequences.  The lineage through Tao, Davenport, Mirsky, and
TPC-137 is recorded with the restricted roles stated above
\cite{Tao2016LogChowla,Davenport1937,Mirsky1948,TPC137}.

The companion certificate enumerates \(576\) finite rows: \(512\) truth
tables, \(27\) three-variable monomials, eight dimension rows, twelve
finite-CRT fixtures, nine landscape rows, and eight analytic-contract rows.
It is a finite reproduction and regression oracle, not the analytic proof.
The immutable closure contains \(117\) Git objects and three ordered remote
logical objects, hence \(120\) logical sources.  Its ordered all-Git digest
is
\nolinkurl{2c187ec15a427ffb0b06a48679f8419be82152fe16ea914c2a86437549117220},
and its logical digest is
\nolinkurl{9315d7c01651ed8b4d94f98c3e4019ad11e28469ee6722903721db280b9f92eb}.
Tao is the only remote analytic provenance: the locator is Theorem 2,
equation (3), printed/PDF page 3, DOI \nolinkurl{10.1017/fmp.2016.6}; the
Cambridge version is CC BY 4.0 but remains nonvendored by project policy.
The Johnston--Yang and Maynard remote objects are closure-only and are
conservatively marked nonredistributable.  All three offline verifiers make
zero network requests by default.  No remote PDF is vendored, and five
inherited payload identities are excluded from both the package members and
the full tree.  RH-393 adds no new remote source.

The theorem does not cover monomials with three or more odd exponents,
unrestricted three-coordinate truth tables, or a generic multishift
safe-table capacity.  It gives no result for growing \(m\), \(q\), shifts,
periodic masks, or \(X\)-dependent coefficients; no effective uniform
rate, ordinary Ces\`aro average, maximum-before-limit, or adaptive capacity
is asserted.  Nothing here supplies an operator model, trace formula, zero
model, or any of Gates A--E.

\section{Limitations}

The fixed-data restriction is qualitative and essential.  The cutoff
argument loses all uniformity as the number of masks, the shifts, or the
period grows.  The first excluded monomial in three variables is exactly a
three-point unsquared M\"obius correlation; neither RH-392 nor Tao's
two-point theorem controls it.  The \(192\)-table corollary is therefore a
genuine but deliberately typed subclass, not evidence that the remaining
\(320\) tables converge or have nonzero limits.  The density landscape is
an arithmetic statement about fixed configurations and does not optimize a
global compatibility problem among several lag constraints.

\section*{Declarations}

\noindent\textbf{Data availability.}
All finite data, source-lock metadata, exact schemas, and verification code
used in this work are included in the accompanying repository package.  No
external payload is redistributed.

\noindent\textbf{Ethics declaration.}
This mathematical study uses no human participants, animals, personal data,
or sensitive datasets; institutional ethics approval was not applicable.

\noindent\textbf{Author contributions (CRediT).}
The RH research program performed conceptualization, formal analysis,
methodology, software, validation, visualization-free presentation, and
writing (original draft, review, and editing).

\noindent\textbf{Conflict of interest.}
The author declares no conflict of interest.

\noindent\textbf{Funding.}
No external funding was received for this work.

\noindent\textbf{AI-use statement.}
AI-assisted tools supported drafting, finite-certificate implementation,
and consistency checks.  Every theorem statement, proof dependency,
citation role, and released byte identity was independently audited; the RH
research program retains responsibility for the content.

\sloppy\raggedright
\bibliographystyle{plainnat}
\bibliography{references}

\end{document}
